\section{Setting}
\label{sec:setting}

\subsection{Theories, regimes and audits}

A \emph{theory package} in the sense of Foundations~I is a carrier together with a lens
that fixes which distinctions are kept, a completion endomap whose fixed points are the
package's objects, and an audit \citep{TsiokosFoundationsI}. The semantic results of
\Cref{sec:strict-join,sec:admission} need much less: a carrier with an observable map, a
join given by maps to the parents, and a step relation on states. The certificate
language of \Cref{sec:certificates} records claims about packages without modelling their
internal structure.

An \emph{admissibility regime} is the declared standard against which an interaction is
judged: which operations are lawful, which sources may be drawn on, which costs must be
paid and what the audit must record. Every existence, obstruction and cost claim below is
relative to a declared regime, and none transfers to another regime without a record of
what the change preserves.

An \emph{audit} is an append-only list of entries, each with a subject, a disposition and
a source location. Append-only matters: a failed, revoked or withdrawn entry remains in
the record, so a later reader can see that something was tried and refused.

\subsection{Certificates and theorems}
\label{sec:setting-kinds}

The catalog's objects are \emph{certificates}: typed records with Boolean and data fields,
usually with an audit. A certificate comes with a validity predicate, usually a conjunction saying
that the right fields are set and, where the record has one, that the audit is non-empty. Many statements about
certificates are therefore immediate: that a valid certificate has a given field set, that
a credit whose definition requires a field is refused when the field is absent, or that a
criterion is equivalent to the definition it restates. The supplement states what each of the
catalog's 36 entries establishes (\Cref{supp:register}).

\paragraph{Status by mathematics.} Each result has one of three statuses, decided by the
mathematics.
\begin{itemize}[nosep]
  \item A \emph{theorem with an instance} is a general statement together with a concrete
        system, written in this paper, in which its hypotheses are proved, so that the
        conclusion is established for that system.
  \item A \emph{general theorem} is a statement whose hypotheses are ordinary data that can
        be checked system by system; where it has classical instances, they rest on
        cited results.
  \item A \emph{certificate specification} fixes what a record must contain and when it is
        valid, together with the consequences and refusals that follow from that
        definition. Like a template, it says nothing about a particular system until the
        record's fields have been established for that system; establishing them is the
        \emph{open obligation} of the specification.
\end{itemize}
Which results are also checked in Lean is a separate question, recorded in the last column
of \Cref{tab:status} and in \Cref{sec:lean}. The instances in this paper are small finite
models; no instance drawn from a particular scientific theory is carried out.

{\footnotesize
\begin{longtable}{@{}>{\raggedright\arraybackslash}p{3.7cm}>{\raggedright\arraybackslash}p{2.1cm}>{\raggedright\arraybackslash}p{4.5cm}>{\raggedright\arraybackslash}p{4.0cm}@{}}
\caption{The results by status.}\label{tab:status}\\
\toprule
\textbf{Result} & \textbf{Status} & \textbf{Instance or open obligation} & \textbf{Checked in Lean} \\
\midrule
\endfirsthead
\toprule
\textbf{Result} & \textbf{Status} & \textbf{Instance or open obligation} & \textbf{Checked in Lean} \\
\midrule
\endhead
\bottomrule
\endfoot
Split-pair criterion (\Cref{thm:split-criterion}) & general theorem & classical; the three-bit join (\Cref{ex:cube}) & the criterion and the finite scan \\
Strictness by split pairs (\Cref{thm:strict-split}) & theorem with an instance & the three-bit join (\Cref{ex:cube}) & the general statement and the example; not the remark on $a\oplus b$ \\
Refinement (\Cref{prop:refinement}, \Cref{rem:refinement-limit}) & general theorem & the identity refinement of the three-bit join shows the converse fails & both \\
Cheap constructions (\Cref{prop:strict-not-cheap}) & general theorem & none needed & relabelling and coarsening \\
Soundness of the strictness certificate (\Cref{thm:bridge}) & theorem with an instance & the three-bit join with $q=\langle p_A,p_B\rangle$ and $s=p_A$ & the general statement and the example \\
Newman's lemma and unique normal forms (\Cref{thm:newman}, \Cref{cor:normal-forms}) & general theorem & classical; termination is needed (\Cref{ex:four-state}) & both, and the four-state example \\
Finite decision (\Cref{prop:finite-peaks}) & general theorem & none needed & reachability, joinability and peak checks \\
Admission order (\Cref{thm:admission}) & theorem with an instance & admission without constraints (\Cref{ex:free-admission}); one disabling guard lies outside the hypothesis (\Cref{ex:disabling}) & all four parts and both examples \\
Enablement without descent (\Cref{prop:enablement-model}) & theorem with an instance & a three-state model & both parts and the positive criterion \\
No free join (\Cref{prop:no-free-join}) & general theorem & a one-event example with the capability supplied & the exclusion and the example \\
Selection (\Cref{prop:selection}) & general theorem & insertion is not a selection step & both parts \\
Bounded search (\Cref{prop:search}) & general theorem & a two-component rendezvous threshold & all parts \\
The 36 catalog entries (\Cref{sec:certificates,sec:separations}) & certificate specifications & open obligation: the record's fields, for each system; the instances are records of the reference world (\Cref{sec:ttw}); two entries are refuted proposals and two are deferred questions & the validity predicates and their consequences, and the finite assays by evaluation \\
\end{longtable}
}

\paragraph{Standing hypotheses.}
Throughout, all records are well typed and carry source locations; no undeclared
operation, source or observer resource is used; and results from earlier papers enter
only through recorded bridges, which name the source and target and the hypotheses kept,
added and lost. These three hypotheses are stated once here and apply to every
certificate statement below.

\paragraph{Finite evidence.}
Many entries are accompanied by an exhaustive enumeration over a declared finite carrier.
Such an enumeration establishes the status of every case inside that carrier and nothing
about cases outside it. Where a statement is universal, its warrant is the proof.

\subsection{Notation}

We write $\tjoinof{\mathcal T_1}{\mathcal T_2}{\mathcal R}$ for a join of theory packages
$\mathcal T_1,\mathcal T_2$ under a regime~$\mathcal R$. The symbol $\tjoin$ distinguishes
this object from Foundations~V's repair join $\repairjoin$, which is a total lattice
operation on quotients of one carrier; a join of theories relates two packages, need not
exist, and when it exists must be paid for. In \Cref{sec:strict-join} a join is modelled
by a carrier $J$ with maps $p_A:J\to A$ and $p_B:J\to B$; we write
$\langle p_A,p_B\rangle:J\to A\times B$ for their pairing. For a relation $\to$ on states,
$\to^{*}$ is its reflexive and transitive closure.
